\chapter{Empirical Analysis and Baseline Results}\label{chap:empirical-analysis-and-baseline-results}

In this chapter, we present an empirical analysis of the multi-currency portfolio optimization problem using historical data. We apply the model-free return representation derived in chapter \ref{chap:model-free-returns} to compute realized returns for the portfolio strategies derived in section \ref{todo}. We evalutate the performance of the various strategies on two investment universes, a large investment universe consisting of 50 biggest stocks in the 14 largest developed markets and a small investment universe consisting of 10 biggest stocks in the same markets. We compare the performance of the optimized portfolios to a benchmark equally weighted portfolio. 


\section{Backtest Results of Expected Shortfall Optimization Approaches}

\begin{table}%[htbp]
  \centering
  \scalebox{0.82}{%
  \begin{tabular}{lcccccccc}
    \toprule
    \multicolumn{8}{c}{\textbf{Panel A: Small Investment Universe (top 10 stocks per economy)}} \\
    \midrule
    \textbf{Metric} & \textbf{Equal Weight} & \textbf{Assets-Full} & \textbf{Assets-Zero} & \textbf{Overlay} & \textbf{Overlay-Reg} & \textbf{Joint} & \textbf{Joint-Reg}\\
    \midrule
    ES (\%) & & \\
    Vol (\%) & & \\
    MD (\%)  & & \\
    SR   & & \\
    SOR  & & \\
    ESR  & & \\
    $\text{Turnover}^a$  & & \\
    $\text{Turnover}^c$  & & \\
    \midrule
    \multicolumn{8}{c}{\textbf{Panel B: Medium Investment Universe (top 50 stocks per economy)}} \\
    \midrule
    \textbf{Metric} & \textbf{Equal Weight} & \textbf{Assets-Full} & \textbf{Assets-Zero} & \textbf{Overlay} & \textbf{Overlay-Reg} & \textbf{Joint} & \textbf{Joint-Reg}\\
    \midrule
    ES (\%) & & \\
    Vol (\%) & & \\
    MD (\%)  & & \\
    SR   & & \\
    SOR  & & \\
    ESR  & & \\
    $\text{Turnover}^a$  & & \\
    $\text{Turnover}^c$  & & \\
    \midrule
    \multicolumn{8}{c}{\textbf{Panel C: Large Investment Universe (top 100 stocks per economy)}} \\
    \midrule
    \textbf{Metric} & \textbf{Equal Weight} & \textbf{Assets-Full} & \textbf{Assets-Zero} & \textbf{Overlay} & \textbf{Overlay-Reg} & \textbf{Joint} & \textbf{Joint-Reg}\\
    \midrule
    ES (\%) & & \\
    Vol (\%) & & \\
    MD (\%)  & & \\
    SR   & & \\
    SOR  & & \\
    ESR  & & \\
    $\text{Turnover}^a$  & & \\
    $\text{Turnover}^c$  & & \\
    \bottomrule
  \end{tabular}%
}
  \caption{Performance metrics for Expected Shortfall optimization approaches on increasing universe sizes}
    \label{tab:es-universe-sizes}
   \vspace{0.5em}
  \begin{minipage}{\textwidth}
  \raggedright\small
    Table \ref{tab:es-universe-sizes} reports out-of-sample performance metrics for the international
  portfolio optimization under expected shortfall, with portfolios rebalanced
  quarterly. The realized out-of-sample statistics shown are Expected Shortfall (ES), volatility (Vol),
  maximum drawdown (MD), the Sharpe Ratio (SR), the Sortino Ratio (SOR), the Expected Shortfall Ratio (ESR), 
  and average asset and currency turnover, denoted Turnover$^a$ and Turnover$^c$, respectively. 
  Each panel corresponds to a different investment universe:
  Panel A covers the small universe, made up of the 10 largest companies by market
  capitalization in each economy; Panel B the medium universe, comprising the 50 largest
  companies per economy; and Panel C the large universe, consisting of the 100 largest
  companies per economy. The sample period runs from June 1990 to June 2023, and the
  first out-of-sample portfolio is formed in June 2000. Details on the optimization approaches and 
    performance metrics are provided in section \ref{todo}.
  \end{minipage}
\end{table}

We can see in table \ref{tab:es-universe-sizes} that the performance of the expected shortfall optimization approaches varies across 
different universe sizes. In the small universe, the optimized portfolios generally outperform the equal weight benchmark in terms of 
risk-adjusted returns, as measured by the Sharpe Ratio and Sortino Ratio. However, as we increase the universe size to medium and large, 
the performance of the optimized portfolios relative to the equal weight benchmark deteriorates, with some strategies even underperforming 
the benchmark in terms of risk-adjusted returns. This suggests that the benefits of optimization may diminish as the investment universe expands, 
potentially due to increased estimation error and model uncertainty in larger universes. 

\section{Backtest Results of Mean-Variance Optimization Approaches}

\begin{table}[htbp]
  \centering
  \scalebox{0.82}{%
  \begin{tabular}{lcccccccc}
    \toprule
    \multicolumn{8}{c}{\textbf{Panel A: Small Investment Universe (top 10 stocks per economy)}} \\
    \midrule
    \textbf{Metric} & \textbf{Equal Weight} & \textbf{Assets-Full} & \textbf{Assets-Zero} & \textbf{Overlay} & \textbf{Overlay-Reg} & \textbf{Joint} & \textbf{Joint-Reg}\\
    \midrule
    ES (\%) & & \\
    Vol (\%) & & \\
    MD (\%)  & & \\
    SR   & & \\
    SOR  & & \\
    ESR  & & \\
    $\text{Turnover}^a$  & & \\
    $\text{Turnover}^c$  & & \\
    \midrule
    \multicolumn{8}{c}{\textbf{Panel B: Medium Investment Universe (top 50 stocks per economy)}} \\
    \midrule
    \textbf{Metric} & \textbf{Equal Weight} & \textbf{Assets-Full} & \textbf{Assets-Zero} & \textbf{Overlay} & \textbf{Overlay-Reg} & \textbf{Joint} & \textbf{Joint-Reg}\\
    \midrule
    ES (\%) & & \\
    Vol (\%) & & \\
    MD (\%)  & & \\
    SR   & & \\
    SOR  & & \\
    ESR  & & \\
    $\text{Turnover}^a$  & & \\
    $\text{Turnover}^c$  & & \\
    \midrule
    \multicolumn{8}{c}{\textbf{Panel C: Large Investment Universe (top 100 stocks per economy)}} \\
    \midrule
    \textbf{Metric} & \textbf{Equal Weight} & \textbf{Assets-Full} & \textbf{Assets-Zero} & \textbf{Overlay} & \textbf{Overlay-Reg} & \textbf{Joint} & \textbf{Joint-Reg}\\
    \midrule
    ES (\%) & & \\
    Vol (\%) & & \\
    MD (\%)  & & \\
    SR   & & \\
    SOR  & & \\
    ESR  & & \\
    $\text{Turnover}^a$  & & \\
    $\text{Turnover}^c$  & & \\
    \bottomrule
  \end{tabular}%
}
  \caption{Performance metrics for mean-variance optimization approaches on increasing universe sizes}
  \label{tab:mv-universe-sizes}
   \vspace{0.5em}
  \begin{minipage}{\textwidth}
  \raggedright\small
    Table \ref{tab:mv-universe-sizes} reports out-of-sample performance metrics for the international
  portfolio optimization under mean-variance, with portfolios rebalanced
  quarterly. The realized out-of-sample statistics shown are Expected Shortfall (ES), volatility (Vol),
  maximum drawdown (MD), the Sharpe Ratio (SR), the Sortino Ratio (SOR), the Expected Shortfall Ratio (ESR), 
  and average asset and currency turnover, denoted Turnover$^a$ and Turnover$^c$, respectively. 
  Each panel corresponds to a different investment universe:
  Panel A covers the small universe, made up of the 10 largest companies by market
  capitalization in each economy; Panel B the medium universe, comprising the 50 largest
  companies per economy; and Panel C the large universe, consisting of the 100 largest
  companies per economy. The sample period runs from June 1990 to June 2023, and the
  first out-of-sample portfolio is formed in June 2000. Details on the optimization approaches and 
    performance metrics are provided in section \ref{todo}.
  \end{minipage}
\end{table}

\section{The Effects of Universe Size on Portfolio Performance}